\hojaejercicio{Cuádricas II}
\ejercicio{Completa las siguientes cuádricas afines de $\mathbb{A}_{\mathbb{C}}^3$ a cuádricas proyectivas en $\mathbb{P}_{\mathbb{C}}^3$ y clasifica las cuádricas obtenidas:
\begin{enumerate}
    \item[a)] $4x^2 + z^2 - 2xy + 2y + 4 = 0$.
    \item[b)] $x^2 + z^2 - xy + 2xz - yz + x + 1 = 0$.
\end{enumerate}}{ 
    
    a) $4x^2 + z^2 - 2xy + 2y + 4 = 0$
    Tenemos que
\[
M_{R_A}(\varphi) = 
\left( \begin{array}{c|ccc}
4 & 0 & 1 & 0 \\
0 & 2 & -1 & 0 \\
1 & -1 & 0 & 0 \\
0 & 0 & 0 & 2
\end{array} \right)
\]

La completación proyectiva es: $4x_1^2 + x_3^2 - 2x_1x_2 + 2x_2x_0 + 4x_0^2 = 0$

Es decir
\[
M_{R_P}(\varphi) = 
\begin{pmatrix}
4 & 0 & 1 & 0 \\
0 & 2 & -1 & 0 \\
1 & -1 & 0 & 0 \\
0 & 0 & 0 & 2
\end{pmatrix}
\]

Las cuádricas complejas se clasifican por rango:

\[
|M_{R_P}(\tilde{\varphi})| = 
\left| \begin{matrix}
4 & 0 & 1 & 0 \\
0 & 2 & -1 & 0 \\
1 & -1 & 0 & 0 \\
0 & 0 & 0 & 2
\end{matrix} \right|
= 2 
\left| \begin{matrix}
4 & 0 & 1 \\
0 & 2 & -1 \\
1 & -1 & 0
\end{matrix} \right|
= 2(-2-4) \neq 0
\]

\[
\text{Rang}(M_{R_P}(\tilde{\varphi})) = 4 \implies \text{es proyectivamente equivalente a}
\]
\[
x_0^2 + x_1^2 + x_2^2 + x_3^2 = 0.
\]

b) Hacemos lo mismo para $x^2+z^2-xy+2xz-yz+x+1=0$

\[
\begin{aligned}
M_{R_P}(\tilde{\varphi}) &= 
\begin{pmatrix}
1 & 1 & 0 & 0 \\
1 & 0 & -1/2 & 1 \\
0 & -1/2 & 0 & -1/2 \\
0 & 1 & -1/2 & 1
\end{pmatrix}
\sim
\begin{pmatrix}
1 & 0 & 0 & 0 \\
0 & -2 & -1/2 & 1 \\
0 & -1/2 & 0 & -1/2 \\
0 & 1 & -1/2 & 0
\end{pmatrix}
\sim
\begin{pmatrix}
1 & 0 & 0 & 0 \\
0 & -2 & 0 & 0 \\
0 & 0 & -1/4 & 1/2 \\
0 & 1 & 1/2 & 0
\end{pmatrix} \\
&\sim
\begin{pmatrix}
1 & 0 & 0 & 0 \\
0 & -2 & 0 & 0 \\
0 & 0 & -1/4 & 1/2 \\
0 & 0 & 1/2 & 1
\end{pmatrix}
\sim
\begin{pmatrix}
1 & 0 & 0 & 0 \\
0 & -2 & 0 & 0 \\
0 & 0 & -1/4 & 0 \\
0 & 0 & 0 & 3
\end{pmatrix}
\sim
\begin{pmatrix}
1 & 0 & 0 & 0 \\
0 & 1 & 0 & 0 \\
0 & 0 & 1 & 0 \\
0 & 0 & 0 & 1
\end{pmatrix}.
\end{aligned}
\]

Vemos de nuevo que es proyectivamente equivalente a $x_0^2 + x_1^2 + \dots x_3^2 = 0$
    

}

%ejercicio2
\ejercicio{Completa las siguientes cuádricas afines de $\mathbb{A}_{\mathbb{R}}^3$ a cuádricas proyectivas en $\mathbb{P}_{\mathbb{R}}^3$ y clasifica las cuádricas obtenidas:
\begin{enumerate}
    \item[a)] $x^2 - y^2 + 4xz + 2yz + 2y - 2z - 1 = 0$.
    \item[b)] $x^2 + y^2 + 4z^2 - 2xy + 4xz - 4yz - 4y = 0$.
    \item[c)] $x^2 - 2z^2 - xz + x + z = 0$.
\end{enumerate}}{ }

%ejercicio3
\ejercicio{Clasifica las siguientes cónicas afines reales:
\begin{enumerate}
    \item[a)] $1 + 2x - 6y + 4x^2 - 6xy + y^2 = 0$.
    \item[b)] $1 + 2x - 6y + 4x^2 - 4xy + y^2 = 0$.
    \item[c)] $1 + 2x - 6y + 4x^2 - 2xy + y^2 = 0$.
\end{enumerate}}{ }

%ejercicio4
\ejercicio{Se consideran en $\mathbb{P}_{\mathbb{R}}^2$ la cónica $\widetilde{\varphi}$ de ecuación $x_0^2 - 2x_0x_2 + 4x_1x_2 = 0$ y las rectas $r_{\alpha}: \alpha x_0 - x_1 + x_2 = 0$, con $\alpha \in \mathbb{R}$. Clasifica las cónicas afines inducidas por $\widetilde{\varphi}$ en los planos afines $\mathbb{P}_{\mathbb{R}}^2 \setminus r_{\alpha}$.}{
4) $P^2_{\mathbb{R}}$ tenemos
\[
M_R(\tilde{\varphi}) = 
\begin{pmatrix} 
1 & 0 & -1 \\ 
0 & 0 & 1 \\ 
-1 & 1 & 0 
\end{pmatrix} 
\sim 
\begin{pmatrix} 
1 & 0 & 0 \\ 
0 & 0 & 1 \\ 
0 & 1 & -2 
\end{pmatrix} 
\sim 
\begin{pmatrix} 
1 & 0 & 0 \\ 
0 & 1 & 0 \\ 
0 & 0 & -2 
\end{pmatrix}
\]

con nuestra recta del infinito:
\[
r_\infty \equiv \alpha x_0 - x_1 + x_2 = 0
\]

Hagamos entonces el cambio de referencia:
\[
\begin{cases}
y_0 = \alpha x_0 - x_1 + x_2 \\
y_1 = x_0 \\
y_2 = x_2
\end{cases}
\Rightarrow
\begin{cases}
-y_0 + \alpha y_1 + y_2 = x_1 \\
y_1 = x_0 \\
y_2 = x_2
\end{cases}
\Rightarrow
N_{R'R} = 
\begin{pmatrix} 
0 & 1 & 0 \\ 
-1 & \alpha & 1 \\ 
0 & 0 & 1 
\end{pmatrix}
\]

Entonces, la cuádrica proyectiva en esa referencia es:
\[
M_{R'}(\tilde{\varphi}) = M_{R'R}^t M_R(\tilde{\varphi}) M_{R'R} = 
\begin{pmatrix} 
0 & 0 & -2 \\ 
0 & 1 & 2\alpha - 1 \\ 
-2 & 2\alpha - 1 & 4 
\end{pmatrix}
\]

Sabemos que nuestra cuádrica de infinito es $C_\infty = \begin{pmatrix} 1 & 2\alpha-1 \\ 2\alpha-1 & 4 \end{pmatrix}$.\\
Y al estar en dimensión 2, las posibles signaturas son:
\begin{itemize}
    \item Si $|C_\infty| = 0 \to \text{Sig}(C_\infty) = (1,0)$
    \item Si $|C_\infty| > 0 \to \text{Sig}(C_\infty) = (2,0)$
    \item Si $|C_\infty| < 0 \to \text{Sig}(C_\infty) = (1,1)$
\end{itemize}

Entonces:
\[
|C_\infty| = 
\left| \begin{matrix} 
1 & 2\alpha-1 \\ 
2\alpha-1 & 4 
\end{matrix} \right| 
= 4 - (2\alpha-1)^2 = 0 \iff (2\alpha-1)^2 = 4
\]
\[
\Rightarrow 
\begin{cases} 
2\alpha-1 = 2 \to \alpha = 3/2 \\ 
2\alpha-1 = -2 \to \alpha = -1/2 
\end{cases}
\]

Si $\alpha < -1/2$ ó $\alpha > 3/2 \Rightarrow |C_\infty| < 0 \Rightarrow \text{Sig}(C_\infty) = (1,1) \Rightarrow C_\infty \sim \begin{pmatrix} 1 & 0 \\ 0 & -1 \end{pmatrix}$.\\
Entonces $C_\infty$ tiene dos puntos $\Rightarrow C_{\varphi_A} \sim X_1^2 - X_2^2 = 1$ hipérbola.

Si $\alpha = 3/2$ ó $\alpha = -1/2 \Rightarrow |C_\infty| = 0 \Rightarrow \text{Sig}(C_\infty) = (1,0) \Rightarrow C_\infty \sim \begin{pmatrix} 1 & 0 \\ 0 & 0 \end{pmatrix}$.\\
$C_\infty$ tiene un punto $\Rightarrow C_{\varphi_A} \sim X_1 = X_2^2$ parábola.

Si $-1/2 < \alpha < 3/2$ entonces $|C_\infty| > 0 \Rightarrow \text{Sig}(C_\infty) = (2,0) \Rightarrow C_\infty \sim \begin{pmatrix} 1 & 0 \\ 0 & 1 \end{pmatrix}$.\\
$C_\infty \sim X_1^2 + X_2^2 = 0$ no tiene ningún punto $\Rightarrow C_{\varphi_A} \sim X_1^2 + X_2^2 = 1$.}

%ejercicio5
\ejercicio{Sea $\overline{\varphi}$ la cuádrica de $\mathbb{P}_{\mathbb{R}}^3$ de ecuación $x_0^2 - x_1^2 + 2x_2x_3 = 0$ y $H$ el plano $x_0 - x_1 = 0$. Clasifica la cuádrica afín definida por $\overline{\varphi}$ en el espacio afín $X = \mathbb{P}_{\mathbb{R}}^3 \setminus H$.}{ 
    $C_{\tilde{\varphi}} \equiv x_0^2 - x_1^2 + 2x_2x_3 = 0$ por lo tanto:

\[
M_R(\tilde{\varphi}^2) =
\begin{pmatrix}
1 & 0 & 0 & 0 \\
0 & 0 & 0 & 0 \\
0 & 0 & -1 & 1 \\
0 & 0 & 1 & 0
\end{pmatrix}
\sim
\begin{pmatrix}
1 & 0 & 0 & 0 \\
0 & -1 & 0 & 0 \\
0 & 0 & 2 & 1 \\
0 & 0 & 1 & 0
\end{pmatrix}
\sim
\begin{pmatrix}
1 & 0 & 0 & 0 \\
0 & -1 & 0 & 0 \\
0 & 0 & 2 & 0 \\
0 & 0 & 0 & -1
\end{pmatrix}
\sim
\begin{pmatrix}
1 & 0 & 0 & 0 \\
0 & 1 & 0 & 0 \\
0 & 0 & -1 & 0 \\
0 & 0 & 0 & -1
\end{pmatrix}
\]

Es decir $\tilde{\varphi}$ es no degenerada de signatura (2,2)

Tenemos entonces que $H_\infty \equiv x_0 - x_1 = 0 \implies -x_1 \neq 0$

Cambio de referencia:

\[
\begin{cases}
y_0 = x_0 - x_1 \\
y_1 = x_1 \\
y_2 = x_2 \\
y_3 = x_3
\end{cases}
\implies
\begin{cases}
x_0 = y_0 + y_1 \\
x_1 = y_1 \\
x_2 = y_2 \\
x_3 = y_3
\end{cases}
\implies
M_{R'R} =
\begin{pmatrix}
1 & 1 & 0 & 0 \\
0 & 1 & 0 & 0 \\
0 & 0 & 1 & 0 \\
0 & 0 & 0 & 1
\end{pmatrix}
\]

\[
M_{R'}(\tilde{\varphi}) = M_{R'R}^t M_R(\tilde{\varphi}) M_{R'R} =
\begin{pmatrix}
1 & 0 & 0 & 0 \\
1 & 1 & 0 & 0 \\
0 & 0 & 1 & 0 \\
0 & 0 & 0 & 1
\end{pmatrix}
\begin{pmatrix}
1 & 0 & 0 & 0 \\
0 & 0 & 0 & 0 \\
0 & 0 & -1 & 1 \\
0 & 0 & 1 & 0
\end{pmatrix}
\begin{pmatrix}
1 & 1 & 0 & 0 \\
0 & 1 & 0 & 0 \\
0 & 0 & 1 & 0 \\
0 & 0 & 0 & 1
\end{pmatrix}
\]

\[
=
\begin{pmatrix}
1 & 0 & 0 & 0 \\
1 & 0 & 0 & 0 \\
0 & 0 & -1 & 1 \\
0 & 0 & 1 & 0
\end{pmatrix}
\begin{pmatrix}
1 & 1 & 0 & 0 \\
0 & 0 & 0 & 0 \\
0 & 0 & -1 & 1 \\
0 & 0 & 1 & 0
\end{pmatrix}
=
\begin{pmatrix}
1 & 0 & 0 & 0 \\
0 & -1 & 0 & 0 \\
0 & 0 & -1 & 1 \\
0 & 0 & 1 & 0
\end{pmatrix}
\]

Por lo que $M_{R'}(\tilde{\varphi}) = M_R(\tilde{\varphi})$

ahora bien, nuestra carta afín es $\{X_0 \neq 0\}$ por lo que nuestra cuádrica de infinito $C_\infty \equiv -X_1^2 + 2X_2X_3 = 0 \sim -X_1^2 + X_2^2 + X_3^2 = 0$

que al ser no singular se concluye que:

$\varphi_A$ es afínmente a un hiperboloide de una hoja,

es decir, $\varphi_A \sim x^2 + y^2 = z^2 + 1$.
}

%ejercicio6
\ejercicio{Clasifica la cuádrica afín real
\[ (3\alpha - 1) + 4\alpha z - x^2 + 4\alpha xy + (3\alpha - 1)y^2 - z^2 = 0 \]
atendiendo a los valores del parámetro $\alpha \in \mathbb{R}$.}{ }

%ejercicio7
\ejercicio{Clasifica, según los valores de $t \in \mathbb{R}$, la cónica afín real
\[ 1 - 2x + 2y + tx^2 + 2xy + ty^2 = 0 \]
Calcula centro y asíntotas en $t = 0$ y $t = 2$ si los hubiera.}{ }

%ejercicio8
\ejercicio{Halla la ecuación de la cónica que es tangente a las rectas:
$x_2 = 0$, $x_0 - x_1 + x_2 = 0$, $x_0 + x_2 = 0$, $x_0 + x_1 + 2x_2 = 0$ y $x_0 + 2x_1 - x_2 = 0$.}{ }

%ejercicio9
\ejercicio{Encuentra todas las cónicas degeneradas del haz generado por
$x_0x_2 = x_1^2$ y $x_0x_1 + 2x_0x_2 = x_1x_2 + 2x_1^2$.}{ }

%ejercicio10
\ejercicio{Obtén la intersección de las cónicas:
\[ x_0^2 + 2x_0x_1 + 2x_0x_2 + x_1x_2 + x_2^2 = 0 \]
\[ x_0^2 + 2x_0x_1 + 3x_0x_2 + 2x_1x_2 + 2x_2^2 = 0 \] }{ }